\documentclass[11pt,a4paper]{article}

\usepackage[margin=1.05in]{geometry}
\usepackage{amsmath,amssymb,amsthm}
\usepackage{microtype}
\usepackage{enumitem}
\usepackage{booktabs}
\usepackage{xcolor}
\usepackage[colorlinks=true,linkcolor=black,urlcolor=teal,citecolor=black]{hyperref}

\definecolor{decision}{RGB}{140,70,10}

\newtheorem*{thmA}{Theorem A (achievability)}
\newtheorem*{thmB}{Theorem B (weak converse)}
\newtheorem*{thmBp}{Theorem B$'$ (weak converse, $\limsup$ form)}

\newenvironment{decisionblock}[1]
  {\medskip\par\noindent\textcolor{decision}{\rule{\textwidth}{0.8pt}}\par\nobreak
   \smallskip\noindent\textbf{\textcolor{decision}{#1}}\par\nobreak\smallskip}
  {\par\noindent\textcolor{decision}{\rule{\textwidth}{0.4pt}}\medskip\par}

\newcommand{\X}{\mathcal{X}}
\newcommand{\Y}{\mathcal{Y}}
\newcommand{\C}{\mathcal{C}}
\newcommand{\Pe}{P_{\mathrm{e}}}
\newcommand{\code}{\mathcal{C}}

\setlist[itemize]{itemsep=2pt,topsep=4pt}

\title{\bfseries The Coding Theorem for Finite\\ Discrete Memoryless Channels\\[4pt]
       \large A formalization submitted for mathematical review}
\author{}
\date{2 August 2026}

\begin{document}
\maketitle

\begin{abstract}
\noindent
Both directions of Shannon's channel coding theorem for finite discrete
memoryless channels --- achievability strictly below capacity, and the weak
converse --- have been formalized and machine-checked, with no unproved
obligations anywhere in the development. This note states what is proved in
ordinary notation, gives the spine of each proof at the level a reviewer needs,
and asks for a decision on the definitional choices the proof now rests on. It
is written so that the mathematics can be evaluated without reading any Lean.

The choice most in need of judgement is how mutual information is defined:
currently by inclusion--exclusion, $I(A;B) = H(A)+H(B)-H(A,B)$, rather than as a
relative entropy. That choice is already load-bearing under a complete proof, so
the request in \S3 is a decision with a cost attached, not a preference poll.
\end{abstract}

\section{What is proved}

\subsection{Setting}

Let $\X$ and $\Y$ be finite alphabets. A \emph{discrete memoryless channel} is a
stochastic matrix $W = \bigl(W(y\mid x)\bigr)_{x\in\X,\,y\in\Y}$, with $n$-fold
extension
\[
  W^{n}(y\mid x) \;=\; \prod_{i=1}^{n} W(y_i \mid x_i),
  \qquad x\in\X^{n},\; y\in\Y^{n}.
\]
For an input distribution $p$ on $\X$ write $\mu_p(x,y) = p(x)W(y\mid x)$ for the
joint law and $(pW)(y) = \sum_{x} p(x)W(y\mid x)$ for the output marginal. The
mutual information $I(p;W)$ is that of $\mu_p$, defined in \S2.1 --- \emph{this is
the definition under review} --- and
\[
  C(W) \;=\; \sup_{p}\, I(p;W).
\]

A \emph{block code} of length $n$ is a triple $\code = (M,f,g)$ with $M \ge 1$
messages, a deterministic encoder $f : \{1,\dots,M\}\to\X^{n}$ and a
deterministic decoder $g : \Y^{n}\to\{1,\dots,M\}$. Messages are uniform, and
\[
  R(\code) = \frac{\log_2 M}{n},
  \qquad
  \Pe(\code) = \frac1M \sum_{m=1}^{M}
    \Pr\bigl[\, g(Y^{n}) \neq m \,\big|\, X^{n} = f(m) \,\bigr].
\]

\subsection{The two theorems}

\begin{thmA}
Let $\X$ be nonempty, let $R < C(W)$ and let $\varepsilon>0$. Then for all
sufficiently large $n$ there exists a block code $\code$ of length $n$ with
\[
  R \;\le\; R(\code)
  \qquad\text{and}\qquad
  \Pe(\code) \;\le\; \varepsilon .
\]
\end{thmA}

\begin{thmB}
Let $(\code_n)_{n\in\mathbb N}$ be block codes, $\code_n$ of length $n$, with
$R \le R(\code_n)$ for every $n$ and $\Pe(\code_n)\to 0$. Then $R \le C(W)$.
\end{thmB}

\begin{thmBp}
If the rates $\bigl(R(\code_n)\bigr)_n$ are bounded above and
$\Pe(\code_n)\to 0$, then $\limsup_{n} R(\code_n) \le C(W)$.
\end{thmBp}

Theorem A carries the hypothesis that $\X$ is nonempty; Theorem B does not need
it. Theorem B$'$ is a separate statement rather than a restatement of B, and its
boundedness hypothesis is the point --- see \textbf{D-10} in \S3.2.

\subsection{Scope, stated plainly}

The result is exactly this and no more.

\begin{itemize}
  \item \textbf{Finite} input and output alphabets. Nothing about continuous or
        countably infinite channels, and nothing about Shannon--Hartley.
  \item \textbf{Average} error probability over uniform messages. \emph{Not}
        maximal error --- that would require an expurgation argument, and none is
        built.
  \item \textbf{Weak} converse: above capacity the error probability cannot
        vanish. It is \emph{not} shown that the error tends to $1$; that is the
        strong converse, and it is out of scope.
  \item Nothing is claimed at $R = C(W)$. Theorem A is strict below capacity,
        Theorem B strict above.
  \item Fixed-length block codes, deterministic encoder and decoder, no feedback.
        Zero-error capacity does not appear.
\end{itemize}

\section{The spine of the proof}

\subsection{Entropy and mutual information}

With the convention $0\log_2 0 := 0$,
\[
  H(p) = -\sum_{a} p(a)\log_2 p(a),
  \qquad
  H(A\mid B) := H(A,B) - H(B),
\]
and, for a joint law on $\mathcal A\times\mathcal B$,
\[
  \boxed{\;I(A;B) \;:=\; H(A) + H(B) - H(A,B).\;}
\]
The boxed definition is \textbf{D-5}, the choice most in need of judgement
(\S3.1).

Two consequences are used throughout. First, the chain rule
\[
  H(A) \;=\; I(A;B) + H(A\mid B)
\]
is an \emph{algebraic identity} under this definition --- substitute and cancel
--- so it holds with no side conditions at all. Second, the entire development
rests on one analytic input, $\log t \le t-1$, from which Gibbs' inequality, the
maximum-entropy bound $H(p)\le\log_2|\mathcal A|$ and subadditivity all follow.
There is no appeal to Jensen's inequality or to convexity anywhere.

\subsection{The converse}

Let $\code$ have $M$ messages and length $n$, with $M$ uniform, $X^{n}=f(M)$ and
$Y^{n}\sim W^{n}(\cdot\mid X^{n})$. Then
\[
\begin{aligned}
  \log_2 M
  \;&\overset{(1)}{=}\; H(M)
  \;\overset{(2)}{=}\; I(M;Y^{n}) + H(M\mid Y^{n})\\[2pt]
  \;&\overset{(3)}{\le}\; I(X^{n};Y^{n}) + H(M\mid Y^{n})
  \;\overset{(4)}{\le}\; n\,C(W) + 1 + \Pe\log_2 M .
\end{aligned}
\]
Here (1) is the entropy of the uniform law, (2) is the chain rule and hence an
identity, (3) is data processing along $M\to X^{n}\to Y^{n}$, and (4) combines
single-letterization with Fano's inequality. Dividing by $n$ gives the single
per-block-length inequality that both asymptotic forms consume,
\[
  R(\code)\bigl(1-\Pe(\code)\bigr) \;\le\; C(W) + \frac1n ,
\]
and letting $n\to\infty$ with $\Pe\to 0$ yields Theorem B.

\paragraph{Step (3) is an equality.}
The message-to-output link is the joint law of the \emph{composed} channel
$m \mapsto W^{n}(\cdot\mid f(m))$, and a change of variables along the encoder
identifies $I(M;Y^{n})$ with $I(X^{n};Y^{n})$ exactly. The data-processing
\emph{inequality} is never needed. This was unexpected and is worth your eye.

\paragraph{Single-letterization.}
For any joint law of the input block,
\[
\begin{aligned}
  I(X^{n};Y^{n})
  &= H(Y^{n}) - \sum_{i} H(Y_i\mid X_i)\\[2pt]
  &\le \sum_i H(Y_i) - \sum_i H(Y_i\mid X_i)
   = \sum_i I(X_i;Y_i)
   \;\le\; n\,C(W).
\end{aligned}
\]
Memorylessness enters exactly once, in the first equality, and it is available
\emph{definitionally}: the $n$-fold channel's law \emph{is} the product
$\prod_i W(y_i\mid x_i)$, so ``given $X^{n}$, the $Y_i$ are independent'' is not a
fact to be established. The only genuine inequality is subadditivity of entropy.
The last step bounds each $I(X_i;Y_i)$ by the supremum $C(W)$, which requires
knowing that supremum is over a bounded set --- see \textbf{D-7} in \S3.2.

\paragraph{Fano's inequality,} in the crude form used here, is
\[
  H(M\mid Y^{n}) \;\le\; 1 + \Pe\,\log_2 M .
\]
It is proved by Gibbs' inequality against the explicit reference weight
$w(m,y) = \nu(y)\,r(m\mid y)$, where $\nu$ is the output marginal and
\[
  r(m\mid y) =
  \begin{cases}
    \tfrac12, & m = g(y),\\[2pt]
    \dfrac{1}{2(M-1)}, & m \neq g(y).
  \end{cases}
\]
Because Gibbs is stated for a \emph{sub-probability} weight ($w\ge 0$,
$\sum w \le 1$) rather than a probability distribution, the degenerate case
$M=1$ needs no separate treatment: the reference has total mass $1$ when
$M\ge 2$ and $\tfrac12$ when $M=1$, and the inequality is indifferent to which.

\subsection{Achievability}

The route is random coding with \emph{threshold decoding on the information
density}, not joint typicality. Define
\[
  i_p(x;y) = \log_2\frac{W(y\mid x)}{(pW)(y)},
  \qquad
  i_p^{\,n}(x;y) = \sum_{j=1}^{n} i_p(x_j;y_j),
\]
so that $\mathbb{E}_{\mu_p}\bigl[i_p(X;Y)\bigr] = I(p;W)$. That identity is
proved, and it is the bridge between the definition in \S2.1 and the quantity the
argument actually concentrates.

\paragraph{(a) A one-shot bound.}
For every threshold $\tau$ and every $M\ge 1$ there exists a code with $M$
messages such that
\[
  \Pe \;\le\;
  \Pr\nolimits_{\mu_p^{\,n}}\bigl[\, i_p^{\,n}(X^{n};Y^{n}) \le \tau \,\bigr]
  \;+\; M\,2^{-\tau}.
\]
The codebook is drawn i.i.d.\ from $p^{\,n}$ and decoded by thresholding the
likelihood ratio. A decoding error means either the true codeword failed the
threshold --- the first term --- or some other codeword passed it, which the
union bound together with a change of measure bounds by $M2^{-\tau}$.

\paragraph{(b) A weak law.}
For $\delta>0$, $\Pr\bigl[i_p^{\,n} \le n(I(p;W)-\delta)\bigr]\to 0$. This is
Chebyshev: on a finite alphabet the information density has finite variance, and
the variance of the i.i.d.\ sum is $n$ times the single-letter variance. No
measure theory is used anywhere in the development; everything is a finite sum.

\paragraph{(c) Assembly.}
Given $R<I(p;W)$, take $\delta = \tfrac12\bigl(I(p;W)-R\bigr)$,
$\tau = n\bigl(I(p;W)-\delta\bigr)$ and $M = \lceil 2^{nR}\rceil$; then (a) and
(b) drive both terms to $0$. For Theorem A one chooses $p$ with $I(p;W)>R$, which
exists because $C(W)$ is a supremum, and the degenerate case $R\le 0$ is handled
by a one-message code.

\paragraph{This is an existence argument, and no encoder falls out of it.}
The final step of (a) is a pigeonhole: since the \emph{average} error over the
ensemble is at most the bound, \emph{some} codebook is at least as good as the
average. That step exhibits a good codebook and gives no procedure for finding
one. It is isolated in a single named lemma precisely so the claim is auditable
rather than folklore.

\paragraph{One technical remark.}
The quantity $i_p(x;y)$ must be $-\infty$ when $W(y\mid x)=0$ and $(pW)(y)>0$.
Under $\mu_p$ this never matters, since that law puts no mass on such pairs; but
the union-bound term integrates against a \emph{product of marginals}, which does
charge them, and there the change-of-measure bound fails outright. The decoder
therefore thresholds on the mass inequality
$2^{\tau}(pW)^{n}(y) < W^{n}(y\mid x)$, in which no logarithm appears, and a
bridging lemma identifies this with $i_p^{\,n}>\tau$ exactly on the support of
the joint law.

\section{What needs your decision}

\noindent
\fbox{\begin{minipage}{0.965\textwidth}
\smallskip
\textbf{The request, in one place.} Decide \textbf{D-5} (how mutual information is
defined). Then a yes or no on \textbf{D-3}, \textbf{D-7}, \textbf{D-8},
\textbf{D-9} and \textbf{D-10}. Each choice below was made in order to get the
proof through; none has been reviewed by anyone with authority over the
mathematics. All are load-bearing under a complete proof, so a rejection has a
cost --- which is stated, so it can be weighed.
\smallskip
\end{minipage}}

\subsection{D-5 --- how mutual information is defined. Decide this one first.}

\begin{decisionblock}{The choice}
\noindent\emph{Current definition:}
\[
  I(A;B) \;=\; H(A) + H(B) - H(A,B).
\]
\noindent\emph{The alternative, and the one most readers would expect:}
\[
  I(A;B) \;=\; \sum_{a,b}\mu(a,b)\,
    \log_2\frac{\mu(a,b)}{\mu_A(a)\,\mu_B(b)}.
\]
\end{decisionblock}

\paragraph{Why the first was chosen.}
It contains no division. The relative-entropy form has a quotient whose
denominator can vanish, so every lemma stated through it carries an implicit
hypothesis that the conditioning event has positive mass --- and that is exactly
the class of hypothesis that gets assumed silently rather than discharged. With
inclusion--exclusion there is nothing to assume: entropy is total, the chain rule
$H(A) = I(A;B)+H(A\mid B)$ is an algebraic identity, and no statement anywhere in
the development needs a support condition.

\paragraph{What it costs.} Two things, both real.
\begin{enumerate}[label=(\arabic*)]
  \item $I(A;B)\ge 0$ \textbf{is not immediate.} Under the relative-entropy form
        it is Gibbs' inequality directly. Here it is subadditivity of entropy,
        which is a theorem --- and in fact it is \textbf{not currently proved},
        because neither of the two theorems needed it. See \S4.
  \item \textbf{Agreement with the relative-entropy form is a lemma, not a
        definition,} and it is not proved. A reader who expects the standard
        formula has to take the identification on trust, or prove it.
\end{enumerate}

\paragraph{The decision.}
If the relative-entropy form should be the definition, say so now. The entropy
toolkit, Fano's inequality, single-letterization and the information-density
bridge all rest on the current choice, so the change is not free --- but it is far
cheaper now than after anything further is built on top. If inclusion--exclusion
is acceptable as the definition, the natural follow-up is to prove the
identification with the relative-entropy form as a lemma, so that the development
connects to the standard literature explicitly rather than by assertion.

\subsection{The other open conventions}

Each is genuinely open. None is merely cosmetic; none is as consequential as
D-5.

\paragraph{D-10 --- the converse is stated over $R$, not $\limsup$, and this is a
soundness issue rather than a stylistic one.}
The natural-looking statement ``$\limsup_n R(\code_n)\le C(W)$ whenever the error
vanishes'' is, as written, \emph{weaker than it appears}. In the formalization
$\limsup$ of a real sequence is an infimum of eventual upper bounds; for a
sequence unbounded above that set is empty, and the infimum evaluates to $0$. An
unbounded rate sequence with vanishing error would therefore satisfy the
$\limsup$ statement vacuously. Theorem B is accordingly stated over a fixed $R$,
and Theorem B$'$ carries an explicit boundedness hypothesis. \emph{Please confirm
that the $R$-form should be primary.}

\paragraph{D-7 --- capacity is a supremum, and attainment is not assumed.}
$C(W)=\sup_p I(p;W)$, with attainment deliberately not a prerequisite for stating
it. The same infimum-of-the-empty-set phenomenon applies: an unbounded set of
achievable mutual informations would make the supremum evaluate to $0$. That the
set is bounded is therefore proved rather than assumed, via
$I \le H(Y) \le \log_2|\Y|$, through the \emph{output} alphabet.

\paragraph{D-3 --- blocks are functions $\{1,\dots,n\}\to\X$, not length-$n$
vectors.}
Mathematically identical; chosen so that products over coordinates are ordinary
finite products. No mathematical content turns on it.

\paragraph{D-8 --- the message count is a natural number $M$, and the size
condition is written once, as $2^{nR}\le M$ over the reals.}
No floor or ceiling appears in any \emph{statement} in the development; ceilings
appear only inside constructions. This was a deliberate guard against the usual
source of silent off-by-one errors in asymptotic rate bookkeeping.

\paragraph{D-9 --- a code does not mention the channel.}
$\code=(M,f,g)$ depends only on the alphabets; the channel enters only when the
error probability is computed. This is the honest reading of what a code is, at
the cost of writing $\Pe(\code,W)$ rather than $\Pe(\code)$.

\section{What is not built}

Stated as plainly as what is. \emph{None of these is a gap in the two theorems},
both of which are complete and unconditional. They are simply not attempted.

\begin{itemize}
  \item \textbf{Attainment of capacity}, $\exists\,p,\ I(p;W)=C(W)$. Not proved.
        Would require compactness of the probability simplex and continuity of
        $I(\cdot\,;W)$. By design this is not a prerequisite for stating $C(W)$
        (D-7).
  \item \textbf{Continuity and concavity of $I(\cdot\,;W)$.} Not proved.
  \item \textbf{Identification of the mutual information of \S2.1 with the
        relative-entropy form.} Not proved. This is the standing cost of D-5.
  \item $\boldsymbol{I(A;B)\ge 0}$. Not proved as a standalone fact --- neither
        theorem needed it. Flagged because a reviewer would reasonably expect it
        to be present.
  \item \textbf{Joint typicality.} No typicality set exists in the development,
        and none is needed: achievability was formalized along the
        information-density route, which requires one concentration estimate
        rather than three. Absent \emph{by choice of route}, not by omission.
  \item Everything in \S1.3 --- maximal error, the strong converse, expurgation,
        feedback, continuous alphabets.
\end{itemize}

\section{Status, and what would help most}

The development is complete: both theorems proved, no unproved obligations, and
the only axioms used are the three standard classical ones on which every result
in the ambient mathematical library already depends. It builds clean.

What it does \emph{not} have is exactly what this note asks for. A proof
assistant can check that the derivations are valid. It cannot check that $C(W)$,
$R(\code)$, $\Pe(\code)$ and $I(p;W)$ mean what an information theorist means by
them. Those are stipulations, and reviewing them is the one part of this that no
machine can do.

The most useful thing to return is a decision on \textbf{D-5}, and a yes or no on
each of D-3, D-7, D-8, D-9 and D-10. If D-5 is to change, it is much cheaper to
change it now than after anything else is built on top.

\vfill
\noindent\rule{\textwidth}{0.4pt}\\[2pt]
{\footnotesize
Sources, and the evidence base from which every claim here is drawn, at
\url{https://github.com/aaslyan/FiniteDMC}. The companion documents are
\texttt{REVIEW.md} (this note in Markdown), \texttt{DOSSIER.md} (evidence, every
claim traced to a build result or a source read), \texttt{GOAL.md} (the twelve
conventions and the obligation ledger) and \texttt{HARD-PARTS.md} (the two hard
obligations, and a junk-value defect found and quarantined while discharging
them).}

\end{document}
